% Chapter 8: Kết luận
\section{Kết luận}

\subsection{Tóm tắt kết quả đạt được}
Đồ án đã hoàn thành đầy đủ các yêu cầu đề ra:

\begin{enumerate}[label=\textbf{\arabic*.}]
    \item \textbf{Agent chơi đúng luật}: Cả ML Agent và Search Agent đều tuân thủ 100\% luật Cờ Tướng, bao gồm các luật đặc biệt (đối mặt tướng, cản chân mã, mắt tượng).

    \item \textbf{Thắng Random 10/10}: ML Agent (Hard) đạt tỷ lệ thắng \textbf{96.7\%} (29/30 ván) với 2 lần đạt 10/10 hoàn hảo. Tất cả chiến thắng đều bằng checkmate.

    \item \textbf{Nhiều mức độ khó}: ML Agent có 3 mức (Easy/Medium/Hard) dựa trên temperature sampling. Search Agent có 10 levels dựa trên depth + heuristic + move ordering.

    \item \textbf{So sánh ML vs Search}: Đã phân tích chi tiết ưu/nhược của cả hai phương pháp. ML nhanh hơn ($\sim$0.5s vs 5--15s) nhưng Search chính xác hơn ở tactical positions.

    \item \textbf{Chỉ dùng Machine Learning}: ML Agent hoàn toàn dựa trên CNN inference + reward shaping, không sử dụng Minimax/Alpha-Beta làm fallback.
\end{enumerate}

\subsection{Đóng góp kỹ thuật}

\subsubsection{Anti-Stalemate Innovation}
Phát hiện quan trọng: ML Agent có xu hướng ăn hết quân đối phương, dẫn đến stalemate (hòa). Giải pháp: kiểm tra terminal state cho top candidate moves, loại bỏ nước gây stalemate và ưu tiên checkmate.

\subsubsection{Canonical Encoding}
Chuẩn hóa góc nhìn (flip bàn cờ cho phe Đen) giúp giảm 50\% pattern cần học, tăng hiệu quả sử dụng dữ liệu.

\subsubsection{Reward Shaping}
Capture bonus phù hợp (Xe=0.8, Pháo=Mã=0.5, Sĩ=Tượng=0.1) giúp agent chơi tactical mà không gây stalemate.

\subsection{Hạn chế}

\begin{itemize}
    \item \textbf{Dữ liệu}: Chỉ huấn luyện trên $\sim$10K samples (v2). Full dataset ($\sim$1M+ samples) sẽ cải thiện đáng kể.
    \item \textbf{Training}: Chưa dùng GPU --- training 5 epochs mất 162s trên CPU.
    \item \textbf{Opening book}: Agent không có tri thức khai cuộc chuyên biệt.
    \item \textbf{Endgame}: Không có endgame tablebase, phụ thuộc vào anti-stalemate heuristics.
    \item \textbf{Self-play}: Chưa implement self-play training (như AlphaZero).
\end{itemize}

\subsection{Hướng phát triển}

\begin{enumerate}
    \item \textbf{Full dataset training}: Huấn luyện trên toàn bộ CCPD dataset ($\sim$120K ván) với GPU.
    \item \textbf{MCTS integration}: Kết hợp Monte Carlo Tree Search với CNN (AlphaZero-style).
    \item \textbf{Self-play}: Cho agent tự chơi với chính mình để cải thiện liên tục.
    \item \textbf{Opening/Endgame databases}: Tích hợp cơ sở dữ liệu khai cuộc và tàn cuộc.
    \item \textbf{Online play}: Triển khai server cho phép đấu trực tuyến.
    \item \textbf{Mobile app}: Port sang nền tảng di động (Flutter/React Native).
\end{enumerate}

\subsection{Bài học kinh nghiệm}

\begin{enumerate}
    \item Pure ML approach có thể đạt hiệu quả cao với reward shaping phù hợp.
    \item Anti-draw mechanisms (stalemate detection, repetition avoidance) là yếu tố quyết định trong trò chơi đối kháng.
    \item Canonical encoding giúp model học hiệu quả hơn đáng kể.
    \item Capture bonus cần được calibrate cẩn thận --- quá cao gây stalemate, quá thấp không đủ aggressive.
\end{enumerate}

\subsection{Phân công công việc nhóm}
\begin{table}[H]
\centering
\caption{Phân công công việc}
\begin{tabularx}{\textwidth}{lX}
\toprule
\textbf{Thành viên} & \textbf{Công việc} \\
\midrule
Dương Đăng Khoa & Data Pipeline: State encoding, tensor representation \\
Nguyễn Đức Gia Bảo & Dataset: Parse PGN/FEN, rule validation \\
Đồng Công Khánh & ML Model: Network architecture, training pipeline \\
Nguyễn Đức Hạnh Nhi & Agent \& Game Loop: ML Agent inference, game logic \\
Trần Hoàng Vỹ Khang & UI, Evaluation \& Report: Pygame UI, benchmarks \\
\bottomrule
\end{tabularx}
\end{table}
